\documentclass[11pt]{article}
\usepackage{amsmath,amsfonts}
\usepackage{graphicx}
%\usepackage{xcolor}
%\definecolor{gray}{rgb}{.75,.75,.75}
%\usepackage[landscape]{geometry}
%\usepackage{hyperref}
%\usepackage[bookmarks=true, pdfborder={0 0 .5}, urlbordercolor=gray]{hyperref}

\renewcommand{\thefootnote}{\fnsymbol{footnote}}


\addtolength{\topmargin}{-.9in}
\addtolength{\textheight}{2in}
\addtolength{\oddsidemargin}{-.4in}
\addtolength{\textwidth}{.8in}
\pagestyle{empty}

\newcommand{\ds}{\displaystyle}
\newcommand{\ts}{\textstyle}

\newcommand{\Z}{\mathbb{Z}}
\newcommand{\R}{\mathbb{R}}
\newcommand{\C}{\mathbb{C}}
\newcommand{\EE}{\mathcal{E}}
\newcommand{\tZ}{\tilde{\Z}}

\newcommand{\Sym}{\mathrm{Sym}}

\renewcommand{\circle}[1]{{\bigcirc\hspace{-.75em}#1}}

\begin{document}

\footnote[0]{Notes by Peter Magyar \ {\tt magyar@math.msu.edu}}
\vspace{-1em}

\centerline{\bf\hfill{\Large\bf
Synthesis: Differential Equations}\hfill }
\vspace{1em}

\noindent 
{\bf Differential equations.} In science and engineering, the solution of differential equations is the main application of calculus, allowing precise quantitative predictions of complicated dynamical behavior.
Here we continue the analysis of exponential growth from \S6.5,
introducing some interesting differential equations 
and solving them by combining an array of techniques from the course. 
\\[.5em]
{\bf Logistic equation.} Consider a process of self-reproduction constrained by an environmental ceiling. That is, the rate of growth of the population $P(t)$ is proportional to the population level, but also to the shortfall between the environmental capacity $E$ (a constant) and the population level: 
$$
\frac{dP}{dt} \ = \ kP\,(E{-}P).
$$
At first, when $P(t)$ stays small and $E{-}P(t)$ stays near $E$, 
we expect $P(t)$ to grow about exponentially;
but as $P(t)$ approaches the capacity $E$,
the shrinking value of $E{-}P(t)$ will slow down the growth, 
making $P(t)$ asymptotically approach $E$.

We solve by separation of variables (\S6.5):
$$
\int \frac{dP}{P(E{-}P)} \ =\ \int k\,dt \ =\ kt + C.
$$
We integrate the left side by a partial fraction decomposition (\S7.4):
$ \frac{1}{P(E{-}P)} \ =\ \frac{A}{P} + \frac{B}{E{-}P}$.
Clearing denominators gives $1=A(E{-}P)+BP$,
 and substituting $P=0$ and $P=E$ allows us to solve for the coefficients:
$A = 1/E$ and $B=1/E$. Thus:
$$
\int \frac{dP}{P(E{-}P)} \ =\
\frac 1E \int \frac{1}{P} + \frac{1}{E{-}P}\ dP 
\ =\ 
\frac 1E(\ln(P) - \ln(E{-}P)) \ =\
\frac 1E\ln\!\left(\frac{P}{E{-}P}\right).
$$
Hence $\frac 1E\ln\!\left(\frac{P}{E{-}P}\right)=kt + C$, which we can solve for $P$ to get a {\it logistic function}:\footnote{
This is closely related to the hyperbolic tangent $\tanh(x) = \frac{e^x-e^{-x}}{e^x+e^{-x}}$
\,(\S6.7).
}
$$
P \ =\ \frac{E
}{1+ Me^{-Ekt}}\ ,
$$
where $M=e^{-EC} = \tfrac{E}{P(0)}-1$\,.
For $E=1$, $k=1$, $P(0)=0.1$, this looks as expected:
\\[1em]
\centerline{
\includegraphics[width=4.5in]
{Synth-3.png}
%\qquad \qquad 
%\includegraphics[width=2in]{1-8-Discont-v.png} 
}
\vspace{0em}

\pagebreak

\noindent
{\sc example:} We model the population of the United States from 1800 to 1900 
by a logistic function fitted to the three census values (in millions of people):
$$
P(1800) = 5.3, \qquad P(1850) = 23.3, \qquad P(1900) = 76.2,
$$
which produces the function:
$$
P(t) \ =\ \frac{182.9}{1 + 33.5\,(0.969)^{t-1800}}\,.
$$
Comparing this graph (red) with the actual decennial census data (black) gives good agreement
even beyond 1900 until the 1950 census. 
\\[1em]
\centerline{
\includegraphics[width=6in]{USPopulationLogistic.png} 
}
\vspace{0em}

\noindent
That is, regardless of territorial expansion, economic booms and depressions, plagues, 
and waves of immigration, the population dynamics from 1800-1940 is remarkably well described 
by the constrained self-reproduction model,
tending toward an eventual population of 183 million.
But the Second World War, which made the United States 
the world's richest and most powerful country
(and introduced antibiotics and vaccination) fundamentally changed US society, 
which is now at 330 million people and growing roughly linearly.
Perhaps this is due to annual immigration quotas?

This is typical of how scientists use calculus to construct and test mathematical models. 
The failure of a model can be just as instructive as its success.

\pagebreak

\noindent
{\bf Catenary curve.} A flexible chain hanging by its own weight forms a familiar curve:
\\[1em]
\centerline{
\includegraphics[width=4.5in]
{Synth-2.png}
%\qquad \qquad 
%\includegraphics[width=2in]{1-8-Discont-v.png} 
}
\\[1em]
To determine the equation $y=f(x)$ of this curve, consider 
a segment stretching from the lowest point $(0,0)$
 to an arbitrary point $(x,y)$. 
 The external forces on this segment are:
a constant horizontal tension $-T_0$ pulling on the left end; gravity pulling
downward with force $-g L(x)$, proportional
to the arclength $L(x)$; and the tangential tension
vector $T(x)$ at the right end, decomposing into components 
$T_1(x)$ and $T_2(x)$.  

With the chain in the equilibrium position,  the above forces must cancel:
$T_1(x)=T_0$ and $T_2(x) = g L(x)$. 
Thus the tangent slope of the curve at $(x,y)$ is:
$
\frac{dy}{dx} = \frac{T_2(x)}{T_1(x)}= k L(x),
$
where $k=g/T_0$.  Applying the formula for arclength of a graph,
(\S8.1), we get:
$$
\frac{dy}{dx} \ =\ k\int_0^x\!\!  \sqrt{\vphantom{|^M}\smash{
1\,{+}\,(\tfrac{dy(t)}{dt})^{2}}}\ dt.
$$
Differentiating both sides by $\frac{d}{dx}$ and using the First Fundamental Theorem (\S4.3) 
gives
$
\frac{d^2y}{dx^2} \ =\ k  \sqrt{\vphantom{m^2}\smash{
1\,{+}\,(\tfrac{dy}{dx})^{2}}},
$
a second order differential equation for $y=f(x)$.
Since only the derivatives of $y$ appear in the equation, not $y$ itself,
we can rewrite it in terms of a new variable $z=\frac{dy}{dx}=f'(x)$, as:
$\frac{dz}{dx} \ =\ k\sqrt{1{+}z^2}$. This is a separable equation (\S6.5):
$$
\int\! \frac{dz}{\sqrt{1{+}z^2}} \ =\ \int k\, dx \ =\  k x + C\,.
$$
The left-hand integral can be determined using the trigonometric substitution (\S7.3)
$z = \tan\theta$, $\sqrt{1{+}z^2}=\sec\theta$, $dz=\sec^2\!\theta\,d\theta$, 
then $\int\!\sec\theta$\, (\S7.2):
$$
\int\! \frac{dz}{\sqrt{1{+}z^2}} \ =\ \int\frac{\sec^2\!\theta}{\sec\theta}\ d\theta
\ =\ \int\sec\theta\ d\theta \ =\ \ln(\tan\theta + \sec\theta)
\ =\ \ln(z+\sqrt{1{+}z^2}).
$$
Solving $\ln(z+\sqrt{1{+}z^2}) = kx + C$ gives: $\frac{dy}{dx}=z = \frac12(e^{kx+C}-e^{-kx-C})$. Integrating:
$$
y \ =\ \frac12\int e^{kx+C}{-}e^{-kx-C}\,dx \ =\ 
\frac{e^{kx+C}+e^{-kx-C}}{2k}+B
\ =\ \frac1k \cosh(kx{+}C) + B,
$$
where we use the hyperbolic cosine $\cosh(x) = \frac12(e^x+e^{-x})$ from \S6.7.
We can choose $k,C,B$ to adjust the curve to any given endpoints and length of chain.
\\[.5em]
{\sc exercise:} Imitate the above arguments 
to work out the shape of a suspension bridge, 
in which the chain has negligeable weight, 
but the gravitational force on a segment 
is proportional to the length of the roadway
suspended beneath: i.e.~proportional to $x$.

\pagebreak

\noindent
{\bf Terminal velocity.} In the standard model of air resistance for a falling object 
or a speeding car, the drag force is proportional to the square of the velocity $v$:
$$
F_d \ =\  \tfrac12 C_d A\rho v^2,
$$
where $C_d$ is  the coefficient of drag, $A$ is cross-sectional area,
and $\rho$ is air density.
In a mass $m$, this force produces the deceleration: 
$$a_d \ =\ \frac{C_d A\rho}{2m} v^2 \ =\ cv^2.$$

If we drop an object as in \S2.7, 
we let $s(t)$ denote its height, $v(t)=s'(t)$ its velocity.
Its total acceleration $a(t)=s''(t)$ combines 
the constant acceleration $-g$ due to gravity with the above drag
deceleration: $a \  =\ -g + cv^2$. This gives the differential equation:
$$
s''(t) \ =\ -g + c(s'(t))^2.
$$
We can simplify this by writing it in terms of $v(t)=s'(t)$ and $\frac{dv}{dt}=v'(t) = s''(t)$:
$$
 \frac{dv}{dt} \ = \ -g+c v^2.
$$
As before, this can be solved by separation of variables (\S6.5):
$$
\int \frac{dv}{g-cv^2} \ =\ \int dt \ =\ -t + C.
$$
The substitution $u=\sqrt{\frac cg}\,v$ leads to an inverse hyperbolic integral (\S6.7):
$$\begin{array}{rcl}
\int \frac{dv}{g-cv^2} &=& 
\frac 1g\int \frac{dv}{1-(\sqrt{\frac cg}\, v)^2} \ =\ 
\frac 1{\sqrt{cg}} \int \frac{du}{1-u^2}
\\[1.3em]
&=&  \frac 1{\sqrt{cg}}\tanh^{-1}(u)
\ = \ \frac 1{\sqrt{cg}}\tanh^{-1}(\sqrt{\frac cg}\,v),
\end{array}$$
where $\tanh^{-1}(u) = \ln\!\sqrt{\frac{1+u}{1-u}}$. Equating this with $-t+C$ and solving for $v$ gives:
$$
 v(t) \ = \ -\sqrt{\tfrac gc}\tanh(\sqrt{gc}\, t +K),
$$
where $\tanh(x) = \frac{e^x-e^{-x}}{e^x+e^{-x}}$. Integrating $v=\tfrac{dy}{dt}$ gives (\S6.7):
$$\begin{array}{rcl}
s(t) &=& -\sqrt{\tfrac gc}\int \tanh(\sqrt{gc}\,t+K)\,dt
\\[.5em]
&=&  -\tfrac1c \ln \cosh(\sqrt{gc}\, t + K) + L,
\\[.5em]
&=&   -\tfrac1c \ln(e^{\sqrt{gc}\, t + K} + e^{-\sqrt{gc}\, t - K}) + M,
\end{array}$$
where $K,M$ are arbitrary constants, and 
$\cosh(x) = \frac{e^x+e^{-x}}{2}$.

If we assume the initial conditions $y(0)=0$, $y'(0)=0$,
this gives $K=0$, $M=\frac 1c\ln(2)$.
Near $t=0$, this has Taylor series (\S11.10) \
$
s(t)=-\tfrac12 g t^2 + \tfrac1{12} c g^2 t^4 +\cdots
$.
Thus at first, for small $t>0$, the falling object closely 
follows the expected ballistic trajectory 
 without air resistance, $s \approx -\frac12 g t^2$.

But eventually, for large $t\to\infty$, we get:
$$
s(t) \ \approx\  -\tfrac1c \ln(e^{\sqrt{c g}\, t}) 
\ =\ -\sqrt{\tfrac gc}\, t \,.
$$
This means the terminal velocity is $v_{\infty}=-\sqrt{\tfrac gc}=-\sqrt{\frac{2mg}{C_d A\rho}}$.
This velocity would be doubled by shrinking cross-section area $A$ by a factor of $\frac14$,
or by dropping 4 times the mass $m$ with the same cross-section.
In the equivalent situation of a car with a constant engine force 
instead of the gravitional force $mg$, 
it means 4 times the horsepower is required to double the top speed.

\newpage

\noindent
{\bf Financial DE: Cumulative inflation.} 
Suppose inflation is recorded by the function $r(t)$,
the annualized inflation rate during a given month $t$.
\\[.5em]
{\sc problem:} What is the cumulative increase in prices up to time $t$? 
\\[.5em] 
Let $y(t)$ be the price level function, the price at time $t$ of an item 
which originally cost \$1 at $t=0$. Counting $t$ in fractions of a year,
denote the time increment by $\Delta t = \frac1{12}$, 
we compute that every month the price undergoes 
inflation of $r(t)(\frac1{12}) = r(t)\Delta t$, giving the new price:
$$
y(t{+}\Delta t) = y(t)\,(1 + r(t)\Delta t) = y(t) + r(t)y(t)\Delta t,
$$
We can rewrite this in terms of $\Delta y=y(t{+}\Delta t)-y(t)$ as:
$$
\frac{\Delta y}{\Delta t} = r(t)y(t).
$$
For given values of $r(t)$, these equations enable us 
to compute the exact monthly amounts $y(t)$ on a spreadsheet.
To get an approximate algebraic solution, 
we imagine the inflation period $\Delta t$ to be
each day, each minute, and eventually the limit $\Delta t\to 0$,
making the finite difference equation into a differential equation 
(which will not change the solution by much):
$$
\frac{dy}{dt} = r(t)y(t).
$$             
Solving this separable DE gives $\int\! \frac{dy}{y} = \int r(t)\,dt$ and:
$$\textstyle
\log y(t) = \int\! r(t)\, dt + C, \qquad 
y(t) = \exp\!\left( \int r(t)\,dt + C\right)  = e^C \exp(\int r(t) \,dt),
$$
and the initial condition $y(0) = 1$ gives:
$$
y(t) = \exp\left(\int_0^t r(x)\,dx\right).
$$
From \S4.4/5.5, recall that the average of a function $f(x)$ over an interval $[a,b]$
is given by the integral $\frac{1}{b-a}\int_a^b f(x)\,dx$. 
Thus the average rate of inflation over $[0,t]$
is 
$$
r_{\text{avg}} = \frac{1}{t}\int_0^t r(x)\,dx,
$$
and $\int_0^t r(x)\,dx   = r_{\text{\!avg}}\, t$. Thus:
$$
y(t) = \exp(r_{\text{\!avg}}\, t).
$$
Note that the average rate of increase $\frac{y(t)-1}t$ is usually 
larger than the average inflation rate  $r_{\text{\!avg}}$, 
because of the effects of compounding:
$$
\frac{y(t)-1}t   \ > \  r_{\text{\!avg}}.
$$
In fact, using the Taylor series for $e^x$, we get
$\frac{y(t)-1}t = r_{\text{\!avg}} + \frac{1}{2!} r_{\text{\!avg}}^2 t
+  \frac{1}{3!} r_{\text{\!avg}}^3 t^2 +\cdots$.
\\[.5em]
{\sc example:} The inflation rate over six months is 
$8.0, 9.0, 8.5, 9.0, 8.5, 8.0$ percent per year, averaging $8.5\%$.
The total price multiple is about 
$y(0.5) = \exp(r_{\!\text{avg}}t) = \exp(0.085\ast 0.5) = 1.0435$,
and the average price rise of $0.0435/0.5 = 0.087 = 8.7\%$ is 
slightly more than the average inflation rate of $8.5\%$.

\newpage

\footnotesize

\noindent
{\bf More Financial DE: Mortgage repayment.} 
Consider a mortgage loan with a fixed annual interest rate $r$, 
and fixed monthly payments at an annual rate of $p$ percent of the loan,
with the current month's interest deducted from each payment. Here $p>r>0$.
\\[.5em]
{\sc problem:} How long to pay off the loan, and how much interest will be paid?
\\[.5em]
Let $y(t)$ be the principal debt at time $t$ years, with $y(0)=1=$ 100\% of the loan,
and denote the montly repayment period by $\Delta t = \frac{1}{12}$. 
Each payment of $p\,\Delta t$ has a deduction 
of the current month's interest $r\,y(t)\Delta t$, so:
$$
y(t{+}\Delta t) \ =\ y(t) - p\,\Delta t + r\,y(t)\Delta t,
$$
which we can rewrite in terms of $\Delta y = y(t{+}\Delta t)- y(t)$ as:
$$
\frac{\Delta y}{\Delta t} \ =\  r\,y(t)- p.
$$
Taking the continuous limit $\Delta t\to 0$:
$$
\frac{dy}{dt} \ =\  r\,y(t)- p.
$$
This is a first order inhomogeneous linear DE (\S6.5) of the form $\frac{dy}{dt} = a(t)y(t)+ b(t)$, 
with solution: 
$y = ce^{A(t)} + e^{A(t)}\!\int\! \frac{b(t)}{e^{A(t)}}\,dt$, 
where $A(t) = \int\! a(t)\,dt = \int r \,dt = rt$ and $b(t) = -p$:
$$
y(t) \ =\  ce^{rt}+ e^{rt}\!\int\! \frac{-p\ }{e^{rt}}\,dt
\ =\  ce^{rt}-p e^{rt}\!\int\! e^{-rt}\,dt
\ =\  ce^{rt} +\frac pr.
$$
Setting $c$ to give the initial condition $y(0)=1$ implies:
$$
y(t) \ =\ \frac pr -\left(\!\frac{p{-}r}{r}\!\right) e^{rt} .
$$
Now we can solve $y(t)=0$ to get the repayment time $T$, and the total interest $I$:
$$
T = \tfrac 1r\ln(\tfrac{p}{p-r}), \qquad\qquad
I = \int_0^T\! r\,y(t)\,dt \ =\ \tfrac pr \ln(\tfrac{p}{p-r}) - 1 \,.
$$
For example, a mortgage at $r = 5\%$ interest with $p=10\%$ annual payment 
takes about $T=20\,\ln(2)= 14$ years to repay, with interest $I =39\%$
of the loan amount.

\medskip

% \newpage

\noindent
{\bf Continuous investment in a fluctuating asset.} 
We regularly invest in an asset whose market price fluctuates according to 
a given function $p(t)$. We want to compute the investment's value $y(t)$ at time $t$.

Over each period from $t$ to $t{+}\Delta t$, 
our previous value fluctuates by a factor of $p(t{+}\Delta t)/p(t)$,
plus an investment of $\Delta t$ (assuming unit annual investment):
$$
y(t{+}\Delta t) \ = \ \frac{p(t{+}\Delta t)}{p(t)} y(t) + \Delta t.
$$
Rewriting in terms of $\Delta y = y(t{+}\Delta t)-y(t)$ and $\Delta p = p(t{+}\Delta t)-p(t)$:
$$
\Delta y + y(t) \ = \ \frac{\Delta p + p(t)}{p(t)} y(t) + \Delta t
\quad\Longrightarrow\quad
\frac{\Delta y}{\Delta t} \ = \ \frac{1}{p}\frac{\Delta p}{\Delta t} y(t) + 1.
$$
In the limit $\Delta t\to 0$, this becomes a first order inhomogeneous linear DE:
$$
\frac{dy}{dt} = \frac 1p \frac{dp}{dt} y + 1 \quad
\text{with} \quad y(0)=0.
$$
By \S6.5, the solution is $y(t) = ce^{A(t)}+e^{A(t)}\!\int \frac{1}{e^{A(t)}}\,dt$
for $A(t) = \int \frac 1p \frac{dp}{dt}\,dt = \ln(p(t))$, so 
$e^{A(t)} = p(t)$. Adjusting for the initial value $y(0)=0$ gives:
$$
y(t) \ = \ p(t)\int_0^t\!  \frac1{p(x)} \,dx.
$$
{\sc examples:}
\begin{itemize}
\item Constant $p(t) = k$ gives $y(t) = t$, exactly equal to our investment,  as expected.

\item A price spike like $p(t) = 1+2t-t^2$ 
with $p(0) = 1$, $p(1)=2$, $p(2)=1$ gives:
$$
y(t) \ =\  
(1{+}2t{-}t^2)\int_0^t\! \frac{1}{2-(x{-}1)^2}\,dx
\ =\ \frac{1{+}2t{-}t^2}{2\sqrt 2}
\ln\!\left(\frac{1+t+\sqrt 2 t}{1+t-\sqrt 2 t}\right),$$
using  $\int\!\! \frac{1}{1-x^2}\,dx = \tanh^{-1}(x) =
\frac12 \ln(\frac{1+x}{1-x})$ (\S6.7)
or partial fractions (\S7.4).

Thus $y(2)\approx 1.25<2$, a loss from our investment: 
the appreciation at the beginning
acts on a smaller asset base than the depreciation at the end.
A price dip would have the opposite effect.
 
%\item A price dip like $p(t) = 2-2t+t^2$ with $p(0) = 2$, $p(1) = 1$, $p(2)=2$ gives:
%$$
%y(t)\ =\  
%(2{-}2t{+}t^2)\int_0^t\! \frac{1}{1+(x{-}1)^2}\,dx
%\ =\ (2{-}2t{+}t^2)(\arctan(t{-}1) +\tfrac 14 {\pi}),
%$$
%an integral using the reverse trig substitution $x{-}1 = \tan\theta$ (\S7.3).
%
%Thus $y(2) = \pi \approx 3.14 > 2$, a profit from our investment,
%reversing the dynamics of the previous example.

\end{itemize}

\end{document}


\noindent   % Moved to Substitution section 
{\bf Medical Math: Heart Flow Rate.} 
In a standard test to measure the rate of blood pumped by the heart, $r$ liters/min, 
doctors inject a colored dye into a vein flowing toward the heart, 
then measure the concentration of dye in arterial blood as it is
pumped out from the heart, $c(t)$ mg/liter after $t$ minutes. 
\\[.5em]
{\sc Problem:} Given the dye concentration function $c(t)$, determine the flow rate $r$.
\\[.5em]
Let the variable $\ell$ denote the liters of blood which have flowed 
through the artery since the start time. 
Assuming the (unknown) flow rate $r$ is constant, we have $\ell=rt$.
Let $C(\ell)$ be the dye concentration after $\ell$ liters have flowed,
so that $C(\ell) = C(rt) = c(t)$.
 
Now, the integral $\int_0^\infty C(\ell)\,d\ell$ sums up:
$$
(\text{mg/liter concentration}) \times (\text{liter increments})
= (\text{mg increments of dye}), 
$$
which computes the total amount of dye, $D$ mg:
$$
D \ =\ \int_0^\infty C(\ell)\,d\ell.
$$
Performing the substitution $\ell=rt$, $d\ell=r\,dt$, we have:
$$
D \ =\ \int_0^\infty C(rt)\,r\,dt \ =\ r\,\int_0^\infty c(t)\,dt.
$$
Then we may compute $r$ as:
$$
r \ = \ \frac{D}{\int_0^\infty c(t)\,dt}.
$$
Since the total dye $D$ is known (the amount injected),
$c(t)$ is measured by the test, and $\int_0^\infty c(t)\,dt$ can
be computed by Riemann sums,\footnote{
Since $c(t)=0$ after all the dye has passed, 
$\int_0^\infty c(t)\,dt$ can be cut off to a finite integral.
}
we obtain flow rate $r$.


%%%%%%%%%%%  Picture  %%%%%%%%%%
\\[0em]
\centerline{
%\includegraphics[width=2in]
{1-8-Discont-iv.png}
\qquad \qquad 
%\includegraphics[width=2in]
{1-8-Discont-v.png} 
}
\vspace{0em}

\noindent
%%%%%%%%%%%%%%%%%%%%%%%%%

%%%%%%%%%%%  Table  %%%%%%%%%%
In fact, we get the following derivatives:
$$\begin{array}{|c||c|c|c|c|c|c|}
\hline &&&&&& \\[-.9em] 
f(x) & \sin(x) & \cos(x) & \tan(x)  & \sec(x) & \csc(x) & \cot(x)
\\[.2em] \hline &&&&&& \\[-.9em]
f'(x) &\cos(x) & -\sin(x) & \sec^2(x) & \tan(x)\sec(x) & -\cot(x)\csc(x) & -\csc^2(x)
\\[.1em] \hline
\end{array}$$
\\[1em]
%%%%%%%%%%%%%%%%%%%%%%%%%%



















